\uuid{KNGZ}
\titre{Fonctions réciproques}
\auteur{bodin}
\organisation{exo7}
\licence{CC BY-NC-SA 4.0}
\datecreate{2020-01-21}
\niveau{L1}
\module{Analyse}
\chapitre{Fonctions circulaires et hyperboliques inverses}
\sousChapitre{Autre}
\difficulte{}
\theme{Réels}
\exo7qcmid{arnaud.arnaud.030}

\contenu{
\begin{question}
\qclassification{126.99}
Quelles sont les assertions vraies ?
\begin{answers}
    \bad{Un fonction continue $\Rr \longrightarrow \Rr$ strictement décroissante est bijective.}
    \good{Si $f$ est une fonction continue bijective croissante alors $f^{-1}$ est croissante.}
    \bad{Si $f$ est une fonction continue bijective ne s'annulant jamais alors $(\frac 1 f)^{-1} = f$.}
    \bad{$\operatorname{arcsin} (\sin x) = x$ pour tout $x \in [0, 2\pi[$.}
    \bad{Si $f(x) = \operatorname{arctan} (x^2)$ alors $f'(x) = \frac{1}{1+x^4}$.}
\end{answers}
\end{question}
}
